\chapter{Marco Metodológico}

El presente proyecto combina tres marcos metodológicos complementarios que
operan en distintas capas del proceso de desarrollo: \textit{Lean Startup}
para la gestión del producto y la validación de valor, el ciclo de vida de
ciencia de datos (CRISP-DM) para el desarrollo técnico del sistema, y
metodología ágil para la organización del trabajo en equipo. La combinación
de estos tres marcos no es accidental: responde a la naturaleza híbrida del
proyecto, que es simultáneamente un producto de software, un sistema de
inteligencia artificial y una herramienta institucional.

\section{Lean Startup y ciclo Build–Measure–Learn}

El enfoque \textit{Lean Startup} se adoptó como marco de gestión del producto,
priorizando la validación empírica de hipótesis sobre la planificación extensa
previa. El principio central aplicado fue: \textit{construir el producto mínimo
que permita aprender, no el producto perfecto}.

El ciclo \textbf{Build--Measure--Learn} se aplicó en cuatro iteraciones
principales del proyecto:

\begin{enumerate}
    \item \textbf{Build:} Prototipo en notebook (Google Colab) con un único
    agente y pipeline básico de segmentación + FAISS.\\
    \textbf{Measure:} 35 inconsistencias detectadas en 47 minutos sobre el
    Contrato A (314 páginas).\\
    \textbf{Learn:} La hipótesis central es técnicamente válida; el cuello de
    botella es la cobertura semántica, no la velocidad.

    \item \textbf{Build:} Arquitectura multi-agente con Jurista, Auditor y
    Cronista en paralelo (v5--v7). Orquestador con \texttt{asyncio.gather}.\\
    \textbf{Measure:} Incremento de cobertura y calidad de hallazgos por
    especialización de prompts.\\
    \textbf{Learn:} Los agentes especializados superan al agente único; los
    esquemas Pydantic son críticos para reproducibilidad.

    \item \textbf{Build:} RAG híbrido (FAISS + BM25 + Cohere Reranker) y
    GraphRAG con extracción de tripletas (v8--v9).\\
    \textbf{Measure:} Mejora en precisión de recuperación normativa y
    enriquecimiento del análisis cruzado entre secciones.\\
    \textbf{Learn:} El grafo más rico no implica más hallazgos (paradoja
    Gemini 3.1); el BM25 es esencial para recuperar cláusulas por número exacto.

    \item \textbf{Build:} Despliegue en producción: Cloud Run, API REST,
    frontend Next.js, bot Telegram, panel de diagnóstico (v9.5--v9.8).\\
    \textbf{Measure:} 119 hallazgos en $\sim$1h 15min sobre Contrato A
    (324 páginas, MVP completo).\\
    \textbf{Learn:} El Cloud Run no soporta SSE nativo; polling cada 5s es
    más robusto. Los \texttt{ContextVar} no propagan en threads de
    \texttt{run\_in\_executor}.
\end{enumerate}

\section{Ciclo de Vida de Ciencia de Datos (CRISP-DM)}

El desarrollo técnico del sistema siguió el proceso estándar de ciencia de
datos CRISP-DM (\textit{Cross-Industry Standard Process for Data Mining}),
adaptado a proyectos de inteligencia artificial generativa:

\begin{table}[H]
\centering
\caption{Aplicación del ciclo CRISP-DM al proyecto ContractIA}
\label{tab:crispdm}
\begin{tabular}{p{4cm}p{9.5cm}}
\toprule
\textbf{Fase CRISP-DM} & \textbf{Aplicación en ContractIA} \\
\midrule
\textbf{1. Comprensión del negocio} &
Entrevistas con 6 especialistas de ProInversión (oct. 2025). Mapeo del
proceso AS-IS. Definición de KPIs SMART. Identificación de la pregunta
central: reducir de 320h a $<$1h. \\
\midrule
\textbf{2. Comprensión de los datos} &
Análisis de 17 contratos de concesión (200--400 páginas). Identificación
de patrones estructurales: numeración jerárquica, referencias cruzadas
densas, tablas de contenido como fuente de falsos positivos. \\
\midrule
\textbf{3. Preparación de los datos} &
Pipeline de extracción (pypdf + pdfminer), normalización UTF-8,
segmentación anti-TOC, indexación multinivel (Secciones, Global,
Local), vectorización con FAISS + BM25 y embeddings
\textit{text-embedding-004}. \\
\midrule
\textbf{4. Modelado} &
Diseño de prompts especializados por agente. Construcción del DiGraph
de conocimiento con NetworkX. Calibración del Cohere Reranker.
Esquemas Pydantic para salida estructurada. Evaluación comparativa
Gemini 2.5 Pro vs.\ Gemini 3.1 Pro Preview. \\
\midrule
\textbf{5. Evaluación} &
Validación sobre contrato real (Contrato A, 324 páginas).
Revisión manual de muestra de 30 hallazgos ALTA por especialistas.
Métricas: precisión 89\%, recall 94\%, cobertura 100\%. \\
\midrule
\textbf{6. Despliegue} &
Containerización Docker. CI/CD con GitHub Actions. Despliegue en
Cloud Run. Interfaces web (Next.js) y conversacional (Telegram).
Disponible en \texttt{contractia.pe}. \\
\bottomrule
\end{tabular}
\end{table}

El ciclo CRISP-DM no se ejecutó de forma lineal sino iterativa: los hallazgos
de la fase de evaluación realimentaron tanto la preparación de datos (mejoras
en el segmentador) como el modelado (refinamiento de prompts), siguiendo la
naturaleza cíclica del proceso.

\section{Desarrollo Ágil}

La gestión operativa del proyecto adoptó prácticas ágiles adaptadas al
contexto de un equipo pequeño (dos integrantes) con restricciones de tiempo
académico:

\begin{itemize}
    \item \textbf{Sprints de dos semanas} con objetivos funcionales claros y
    entregables demostrables al finalizar cada sprint.
    \item \textbf{Versionado semántico} del sistema (v1.0 a v9.8.0) que
    permite trazabilidad precisa de cada funcionalidad incorporada.
    \item \textbf{Integración continua (CI/CD)} mediante GitHub Actions con
    despliegue automático en Google Cloud Run al fusionar cambios en la rama
    principal.
    \item \textbf{Feature flags por variable de entorno} para activar
    funcionalidades experimentales (RAG agéntico, modelos alternativos) sin
    afectar el comportamiento base del sistema.
\end{itemize}

\section{Justificación de la Combinación Metodológica}

La elección de esta combinación responde a tres necesidades concretas del
proyecto:

\begin{enumerate}
    \item \textbf{Lean Startup} aborda la incertidumbre de valor: no se sabía
    a priori si un sistema basado en LLMs podría procesar contratos de 300+
    páginas con suficiente precisión para ser útil. El prototipo temprano
    permitió validar esta hipótesis antes de invertir en infraestructura.
    \item \textbf{CRISP-DM} estructura el trabajo técnico en un dominio de
    alta especialización (Legal NLP) donde la calidad de los datos y la
    pertinencia del modelado son determinantes del éxito. Sin este marco,
    el riesgo de construir un sistema técnicamente correcto pero semánticamente
    débil habría sido elevado.
    \item \textbf{Ágil} permite absorber los aprendizajes de cada ciclo
    Build--Measure--Learn de forma ordenada, sin bloquear el desarrollo por
    planificación excesiva en un dominio donde los requisitos evolucionan con
    el conocimiento del problema.
\end{enumerate}